\documentclass[a4paper]{article}
% Español (México) con XeLaTeX: fontspec para fuentes Unicode, polyglossia
% para la división silábica y los nombres del documento en la variante mexicana.
\usepackage{fontspec}
\usepackage{polyglossia}
\setdefaultlanguage[variant=mexican]{spanish}
% Los nombres chinos van como «pinyin (original)»: xeCJK da a los caracteres
% Han su propia fuente; sin ella XeLaTeX los omite con un aviso.
% FandolSong no cubre algunos caracteres de nombres (U+73FA); WenQuanYi Zen Hei sí.
\usepackage[AutoFallBack=true]{xeCJK}
\setCJKmainfont{FandolSong-Regular.otf}
\setCJKfallbackfamilyfont{\CJKrmdefault}{WenQuanYi Zen Hei}
% polyglossia no traduce los nombres de listings.
\AtBeginDocument{\providecommand{\lstlistingname}{}\renewcommand{\lstlistingname}{Listado}%
  \providecommand{\lstlistlistingname}{}\renewcommand{\lstlistlistingname}{Índice de listados}}
\usepackage{amsmath, amssymb}
\usepackage{graphicx}
\usepackage[margin=2.35cm]{geometry}
\usepackage[most]{tcolorbox}
\usepackage{booktabs}
\usepackage{tabularx}
\usepackage{array}
\usepackage{float}
\usepackage{xcolor}
\usepackage{hyperref}
\usepackage{fancyhdr}
\setCJKmainfont[AutoFakeBold=2.4]{AR PL UMing CN}
\setCJKsansfont[AutoFakeBold=2.4]{Droid Sans Fallback}
\setCJKmonofont[AutoFakeBold=2.4]{Droid Sans Fallback}
\hypersetup{colorlinks=true, linkcolor=blue!55!black, urlcolor=blue!60!black}
\graphicspath{{./}{figures/}}
\setlength{\parskip}{0.55em}
\setlength{\parindent}{2em}
\setlength{\emergencystretch}{3em}
\renewcommand{\arraystretch}{1.18}
\newtcolorbox{knowledgebox}[1]{enhanced, colback=blue!5!white, colframe=blue!70!black, colbacktitle=blue!70!black, coltitle=white, fonttitle=\bfseries, title={#1}, sharp corners, breakable}
\newtcolorbox{importantbox}[1]{enhanced, colback=yellow!10!white, colframe=yellow!75!black, colbacktitle=yellow!75!black, coltitle=black, fonttitle=\bfseries, title={#1}, sharp corners, breakable}
\newtcolorbox{warningbox}[1]{enhanced, colback=red!5!white, colframe=red!70!black, colbacktitle=red!70!black, coltitle=white, fonttitle=\bfseries, title={#1}, sharp corners, breakable}
\pagestyle{fancy}
\fancyhf{}
\fancyfoot[C]{\thepage}
\fancyfoot[R]{\footnotesize\color{gray}\href{https://github.com/hqhq1025/ai-course-notes}{AI Course Notes}}
\renewcommand{\headrulewidth}{0pt}
\renewcommand{\footrulewidth}{0pt}
\newcommand{\videourl}{https://www.youtube.com/watch?v=pWY0HVUH8GA}
\newcommand{\repourl}{https://github.com/hqhq1025/ai-course-notes}

\begin{document}
\begin{titlepage}
\centering
{\Large Notas didácticas de una entrevista a fondo\par}
\vspace{1.0cm}
{\huge\bfseries La escasez de datos en robótica: simulación, datos sintéticos y la cadena industrial de la inteligencia corporizada}\par
\vspace{0.5cm}
{\Large Zhang Xiaojun conversa con Xie Chen (谢晨), fundador de Guanglun Zhineng (光轮智能)\par}
\vspace{0.35cm}
{\large Elaborado a partir del video público de Zhang Xiaojun Podcast, la transcripción por GPU, la estructura de capítulos y diagramas conceptuales propios\par}
\vspace{0.3cm}
{\large 2025-07-15\par}
\vspace{0.85cm}
\includegraphics[width=0.88\textwidth,height=0.42\textheight,keepaspectratio]{cover.jpg}\par
\vfill
\begin{tcolorbox}[width=0.92\textwidth, colback=black!2!white, colframe=black!55, sharp corners]
\textbf{Título del video}: 109. ¿La robótica enfrenta escasez de datos? Conversación con Xie Chen: simulación y datos sintéticos, la adquisición multimillonaria de Meta y Alexandr Wang\par
\textbf{Canal del video}: Zhang Xiaojun Podcast\par
\textbf{Duración del video}: 01:41:09\par
\textbf{Enlace del video}: \href{\videourl}{\nolinkurl{\videourl}}\par
\textbf{Notas sobre el material}: YouTube no tiene subtítulos disponibles; el texto principal se elaboró a partir de la transcripción con faster-whisper large-v3 en CUDA, los capítulos del video, la descripción del video, la portada y figuras didácticas propias. El video de título fijo no se reutiliza en el texto principal.\par
\textbf{Página del proyecto}: \href{\repourl}{\nolinkurl{\repourl}}\par
\end{tcolorbox}
\end{titlepage}

\tableofcontents
\newpage

\section{Introducción: en qué consiste realmente la escasez de datos en robótica}

Esta sección establece primero cómo leer el episodio completo. EP109 no se limita a afirmar de forma general que «los robots necesitan datos», sino que descompone la escasez de datos en cuatro niveles: no hay suficientes robots reales, la interacción física es difícil de recopilar, existe un gap entre la simulación y la realidad y, además, que los datos mejoren de verdad al modelo debe verificarse con un sistema de evaluación. La trayectoria de Xie Chen (谢晨) abarca Cruise, NVIDIA, NIO (蔚来) y Guanglun Zhineng (光轮智能), por lo que su exposición tiene un rasgo muy práctico: cada vez que habla de un concepto técnico, vuelve a preguntarse «si mejora el modelo, si permite llevar el sistema a producción y si cierra un ciclo comercial».

Esta entrevista también aporta una pieza fundamental a la serie sobre robótica de Zhang Xiaojun (张小珺). EP121 aborda, desde la perspectiva de DeepMind, los modelos base para robótica, la generalización entre distintos cuerpos robóticos y los modelos del mundo; EP132 trata, desde la perspectiva de Xinghaitu (星海图), el robot completo, la cadena de suministro y la Data Recipe; EP109, por su parte, explica a fondo el sistema de producción de datos de la base: por qué los datos sintéticos no son «un atajo para sustituir los datos reales», sino una infraestructura que debe calibrarse continuamente con la retroalimentación real, el sistema de evaluación, los parámetros físicos y las necesidades de los clientes.

\begin{importantbox}{Tesis central del episodio}
La escasez de datos de la inteligencia corporizada no se reduce a que «haya muy pocos videos» o «muy pocos robots»; lo que falta son datos del mundo físico que puedan escalarse, verificarse y servir al aprendizaje por refuerzo y al despliegue real. El valor de los datos sintéticos y de la simulación depende de que formen un ciclo cerrado entre Real2Sim, Simulation, Sim2Real y el sistema de evaluación.
\end{importantbox}

El ponente subraya que un mundo simulado de apariencia espectacular que no mejora los modelos de percepción, los modelos de predicción ni las políticas de los robots es solo un juguete vistoso. Una verdadera infraestructura de datos debe diseñarse a partir de los objetivos de aprendizaje automático: qué escenarios son escasos, qué variables influyen en el modelo, qué muestras sintéticas producen una mejora real en la metric y qué gap debe corregirse con recopilación física y retroalimentación del despliegue real.

\begin{figure}[H]
\centering
\includegraphics[width=0.96\textwidth]{sim2real-loop.png}
\caption{Ciclo cerrado Sim2Real: se generan datos en simulación, se despliega en la realidad y se corrige la simulación con la retroalimentación real. Diagrama conceptual propio, elaborado a partir de la conversación entre 00:02:48 y 00:16:17.}
\end{figure}

\begin{knowledgebox}{Lectura de la figura: los datos sintéticos no sustituyen la realidad, la amplifican}
El punto de partida de esta figura no es «generar de la nada», sino los problemas escasos del mundo real: un peatón que aparece de repente, un día con niebla, una pendiente, un agarre fallido, la bisagra de la puerta de un refrigerador, etc. La simulación parametriza, combina y escala estos problemas; después del entrenamiento, se vuelve al mundo real para verificar. Al leer la figura hay que fijarse en las flechas del ciclo cerrado: si la retroalimentación real no regresa a la simulación, los datos sintéticos se vuelven cada vez más un ejercicio de autocomplacencia interna; si el sistema de evaluación no puede medir si el modelo mejora, por grande que sea el volumen de datos no se puede demostrar que sean eficaces.
\end{knowledgebox}

\begin{knowledgebox}{Nota sobre la estrategia visual}
Este video es una toma fija de entrevista, sin diapositivas, pizarrón ni demostraciones de producto. El texto sigue el flujo de trabajo para pódcast: la portada aparece solo en la primera página y el cuerpo no reutiliza imágenes de la entrevista; todas las imágenes del cuerpo son diagramas conceptuales propios que sirven para presentar estructuras didácticas como Sim2Real, Real2Sim, la cadena industrial y la scaling law.
\end{knowledgebox}

\subsection{Resumen de la sección}

La línea principal de EP109 puede condensarse en una frase: para que la inteligencia corporizada pueda hacer scale, no basta con esperar a que los robots reales generen datos de forma natural; hay que construir una infraestructura de datos sintéticos y de simulación que pueda calibrarse con la realidad, evaluarse con los modelos y servir al despliegue real.
\section{Sim2Real, Real2Sim y datos sintéticos}

Esta sección explica primero los términos más frecuentes, porque los juicios posteriores sobre la industria se apoyan en estos conceptos. Sim2Real es ir de la simulación a la realidad; Real2Sim es trasladar la realidad a la simulación; Synthetic Data son datos aptos para entrenamiento generados mediante simulación, algoritmos y procesos programáticos. Los tres no son buzzwords paralelas, sino posiciones distintas dentro de un mismo ciclo de ingeniería: el mundo real plantea el problema, la simulación lo amplía, el modelo se entrena con datos sintéticos y los resultados del despliegue vuelven a calibrar la simulación.

La experiencia de Xie Chen (谢晨) en Cruise aporta una lección decisiva. Al principio, el equipo recurrió a ingenieros de videojuegos para recrear San Francisco; el resultado parecía Matrix, pero el área de percepción no podía usar esos datos directamente, porque no se habían definido según objetivos de aprendizaje automático. Lo que después funcionó de verdad fue construir primero un sistema de evaluación que determinara si los datos sintéticos mejoraban el modelo y, a partir de ello, modificar los recursos, los escenarios, el ruido de los sensores y la cadena de generación de datos.

\begin{knowledgebox}{Términos clave: simulación, datos sintéticos y transferencia}
\begin{tabularx}{\textwidth}{p{0.20\textwidth}p{0.35\textwidth}X}
\toprule
Término & Significado básico & Cuestión central en este episodio \\
\midrule
Sim2Real & Entrenar o generar datos en simulation y luego transferirlos al mundo real & El éxito del despliegue real determina si la simulación sirve. \\
Sim2Real gap & Diferencia de distribución entre la simulación y la realidad & No se busca una diferencia nula, sino determinar a partir de qué magnitud de diferencia los datos empiezan a ser útiles. \\
Synthetic Data & Datos generados mediante simulación, algoritmos, procesos programáticos o colaboración humano-máquina & Lo decisivo no es la cantidad generada, sino la densidad de información, la calidad y la utilidad para el modelo. \\
Real2Sim & Trasladar a la simulación escenarios reales, recursos, parámetros físicos y el propio cuerpo del robot & La reconstrucción visual es relativamente sencilla; lo difícil son los parámetros físicos y la mecánica de la interacción. \\
Evaluation Loop & Inferir la calidad de los datos a partir del entrenamiento del modelo y de los resultados del despliegue real & Sin un sistema de evaluación, al equipo de datos le resulta difícil saber qué corregir. \\
\bottomrule
\end{tabularx}
\end{knowledgebox}

\subsection{De un juguete vistoso a datos para aprendizaje automático}

Esta subsección responde una pregunta concreta: por qué una simulación que «parece real» no equivale a una simulación «útil para el modelo». El sistema de simulación inicial de Cruise podía recrear manzanas enteras de la ciudad, pero si el ruido de los sensores, la distribución de la iluminación, las anotaciones, la probabilidad de los escenarios y los eventos de cola larga no correspondían a lo que requería el entrenamiento del modelo, el sistema de percepción no obtenía ningún beneficio. El primer paso de Xie Chen no fue mejorar la calidad gráfica, sino construir un sistema de evaluación.

El sistema de evaluación tiene dos tipos de indicadores. El primero es el realismo absoluto: si la luz, el color, el ruido de los sensores, la geometría, las anotaciones y la distribución de escenarios coinciden con los datos reales. El segundo es la utilidad: si, tras alimentar al modelo con este conjunto de datos sintéticos, su desempeño en la tarea objetivo mejora. El ponente lo compara con un espresso frente a un café helado rebajado con agua: con el tiempo, los datos se valoran cada vez más por «cuánta información adicional aporta al modelo cada unidad de datos», y no solo por su volumen.

\begin{importantbox}{El sistema de evaluación va primero}
Lo primero que debe preguntarse un equipo de datos sintéticos no es «cómo generar más», sino «qué datos hacen que el modelo mejore». Sin un sistema de evaluación, el equipo de simulación solo puede optimizar la apariencia visual; con él, el equipo puede saber si debe corregir el ruido de los sensores, la distribución de escenarios, los parámetros físicos de los recursos o la cadena de control de calidad de los datos.
\end{importantbox}

La utilidad de los datos puede expresarse como un objetivo simplificado:
\[
U(D_s)=\Delta M_{\mathrm{target}}-\lambda C(D_s)-\gamma G(D_s,D_r)
\]
donde $D_s$ representa los datos sintéticos, $D_r$ los datos reales, $\Delta M_{\mathrm{target}}$ la mejora en el indicador del modelo objetivo, $C(D_s)$ el costo de generación y control de calidad, $G(D_s,D_r)$ el gap relevante entre lo sintético y lo real, y $\lambda$ y $\gamma$ son los pesos del costo y del gap. Esta fórmula no procede textualmente de la entrevista; es una formulación didáctica de su metodología: los datos sintéticos eficaces deben considerar a la vez el beneficio para el modelo, el costo de producción y la diferencia con la realidad.

\subsection{Flujo de datos sintéticos y proporciones}

A continuación, esta subsección aterriza el flujo en un ejemplo: un peatón que aparece de improviso en la calle. Este tipo de corner case es escaso en los datos reales, pero crucial para la seguridad. El flujo de datos sintéticos no consiste en copiar un video diez mil veces, sino en generar variantes sistemáticas en torno a un problema real: distintas edades, complexiones, ropa, trayectorias, carriles, clima, iluminación, condiciones de tráfico y ruido de los sensores. Después, el algoritmo se ejecuta diez mil veces en la simulación para descubrir si la percepción, la predicción y la planificación siguen fallando.

No existe una proporción fija. En el conjunto total de datos de la conducción autónoma, la flota real genera a diario una gran cantidad de datos reales, de modo que los datos sintéticos quizá no sean la mayor parte del total; pero para los corner cases escasos, las muestras reales pueden no bastar ni en un año, por lo que a un solo problema real se le asocian miles o decenas de miles de variantes sintéticas. La experiencia que se comparte en el programa es que, en conjunto, en algún momento cerca del 30\% de los datos fueron sintéticos, mientras que en los escenarios de cola larga la proporción real/sintético puede superar 1:99 o 1:100.

\begin{knowledgebox}{El flujo paso a paso: de un problema a un conjunto de datos de entrenamiento}
\begin{tabularx}{\textwidth}{p{0.20\textwidth}p{0.35\textwidth}X}
\toprule
Paso & Qué se hace & Por qué importa \\
\midrule
Localización del problema real & Encontrar un caso de peligro o de falla & La generación de datos debe partir de riesgos reales. \\
Parametrización del escenario & Variar peatones, carriles, clima, iluminación y flujo de tráfico & Permite que el modelo vea múltiples manifestaciones de un mismo problema. \\
Simulación de sensores & Simular cámaras, radares, ruido y perspectivas de observación & Lo que el modelo recibe al final son datos de sensores. \\
Anotación automática y control de calidad & Generar etiquetas 2D/3D/4D y verificar su calidad & Los errores de anotación contaminan directamente el entrenamiento. \\
Retroalimentación al modelo & Usar los datos para entrenar o evaluar percepción/predicción/planificación & La utilidad debe demostrarse con los indicadores del modelo y los resultados del despliegue. \\
\bottomrule
\end{tabularx}
\end{knowledgebox}

\begin{warningbox}{La proporción no puede discutirse al margen de la tarea}
Decir «1\% de datos reales y 99\% de datos sintéticos» se presta a confusión. La pregunta correcta es: ¿para qué tarea, qué modo de falla, qué módulo del modelo y qué etapa del despliegue tienen los datos sintéticos una utilidad marginal mayor? Las proporciones del conjunto total de datos, de los escenarios de cola larga y de los entornos de aprendizaje por refuerzo pueden ser completamente distintas.
\end{warningbox}

\subsection{Resumen de la sección}

Sim2Real y Real2Sim forman un ciclo cerrado de datos. El valor de los datos sintéticos no lo determina el realismo de la imagen, sino dos tipos de evaluación: si coinciden con las variables decisivas del mundo real y si hacen que el modelo objetivo mejore en tareas reales.
\section{Conducción inteligente vs. inteligencia corporizada: juego visual e interacción física}

La sección anterior explicó cómo funcionan los datos sintéticos en la conducción autónoma; esta sección aborda por qué la inteligencia corporizada es más difícil. La conducción autónoma no es un problema sencillo, pero tiene una enorme ventaja: el automóvil como plataforma existe desde hace más de cien años, una gran cantidad de vehículos circula por caminos reales y las flotas L2+ pueden enviar de forma continua datos reales de visión, de sensores y de conducción. En cambio, no hay millones de robots humanoides o cuadrúpedos funcionando de manera natural en hogares, fábricas y hoteles, por lo que es muy difícil poner en marcha primero el ciclo virtuoso de datos reales.

Por eso Xie Chen (谢晨) sostiene que «la inteligencia corporizada debe apoyarse primero en una gran cantidad de datos sintéticos y después calibrarse con una pequeña cantidad de datos reales». La conducción autónoma se parece más a un juego visual y de reglas: los caminos, los carriles, los participantes del tránsito y el propio vehículo están relativamente bien definidos. La inteligencia corporizada tiene que lidiar con tazas, popotes, refrigeradores, teclados, telas, líquidos, fricción, colisiones y distintos cuerpos de robot. La complejidad de la interacción física impide que dependa solo de las imágenes y los videos que ya existen en internet.

\begin{figure}[H]
\centering
\includegraphics[width=0.94\textwidth]{av-vs-embodied-data.png}
\caption{Conducción inteligente vs. inteligencia corporizada: la conducción inteligente se inclina hacia el juego visual; la inteligencia corporizada depende más de la interacción física. Diagrama conceptual propio, elaborado a partir de la conversación de 00:16:17--00:32:41.}
\end{figure}

\begin{knowledgebox}{Lectura de la figura: por qué el ciclo virtuoso de datos reales es más fácil en la conducción autónoma}
A la izquierda, la conducción autónoma cuenta con una gran cantidad de vehículos reales, un cuerpo estable y una estructura vial, por lo que los datos reales son la mayor parte y los datos sintéticos sirven principalmente para ampliar la cola larga y los problemas de seguridad. A la derecha, la inteligencia corporizada carece de una población de robots que funcionen de manera natural, y sus tareas abarcan objetos, espacio, contacto y acciones; por eso los datos sintéticos y la simulación deben asumir desde el principio el papel de «construir el mundo».
\end{knowledgebox}

\subsection{El cuello de botella de la interacción física}

Esta subsección explica con más detalle qué dificultad añade exactamente la «interacción física». Los datos visuales abundan en internet, y las fotos y los videos pueden ayudar a generar calles, vehículos, árboles y pavimento; pero los datos de interacción física no son solo apariencia, sino fuerza, fricción, ejes de giro, colisiones, materiales, peso, deformación y recuperación tras un fallo. Sin datos de interacción de bajo nivel, es muy difícil que la IA genere de la nada un comportamiento físico confiable.

La puerta del refrigerador es un buen ejemplo de la entrevista. Para que un robot aprenda a abrir un refrigerador, el recurso de simulación no puede ser solo la carcasa de un refrigerador. Necesita el eje de giro de la puerta, las bisagras, la fuerza del cierre magnético, la fuerza necesaria para abrir la puerta en distintos ángulos, los cuerpos de colisión de los cajones y de la puerta, la distribución de distintos modelos de refrigerador y la fuerza de contacto entre la mano del robot y la puerta. No basta con que visualmente parezca un refrigerador: lo decisivo es que físicamente se pueda interactuar con él.

\begin{figure}[H]
\centering
\includegraphics[width=0.96\textwidth]{physical-interaction-bottleneck.png}
\caption{El cuello de botella de la interacción física: el agarre, el contacto, la retroalimentación de fuerza y las diferencias de material hacen que los datos sean especialmente caros. Diagrama conceptual propio, elaborado a partir de la conversación de 00:16:17--00:32:41.}
\end{figure}

\begin{knowledgebox}{Lectura de la figura: la unidad de los datos corporizados no es el frame, sino la interaction}
Cada variable de la figura puede provocar que la tarea falle: si el punto de agarre se desvía un poco, el objeto se resbala; un material distinto cambia la fricción; los errores del actuador se acumulan; y un cuerpo de colisión impreciso hace que la política entrenada en simulación falle en el mundo real. Para la inteligencia corporizada, si un dato no incluye la acción, el resultado mecánico y la retroalimentación del fallo, es muy difícil que le enseñe al robot a realizar tareas de verdad.
\end{knowledgebox}

\begin{warningbox}{No hay que confundir los datos de video con datos de robot}
Los videos de internet pueden mostrarle al modelo «cómo lo hacen los humanos», pero por lo general no contienen el cuerpo del robot, el estado de las articulaciones, las fuerzas de contacto, los comandos de acción ni las variables controlables del entorno. Sirven para el preentrenamiento de conocimiento del mundo, pero no pueden sustituir directamente los datos de interacción en lazo cerrado que requiere el entrenamiento de robots.
\end{warningbox}

\subsection{La vía de la teleoperación y el ciclo comercial}

Esta subsección responde a una objeción frecuente: si los datos reales son importantes, ¿no se podría primero ganar dinero con robots teleoperados y después recolectar datos? Xie Chen considera que esto podría funcionar en algunos escenarios, pero que como vía general difícilmente avanzaría más rápido que una vía basada principalmente en datos sintéticos. La razón es que la teleoperación primero tiene que generar un valor por el que el cliente esté dispuesto a pagar; de lo contrario, solo es un muestreo de alto costo. Además, la teleoperación entre estados o entre países enfrenta regulación, costos de personal, sindicatos, renta de espacios y complejidad operativa.

El ciclo virtuoso de datos de la conducción autónoma es sólido porque el automóvil ya se vende por sí mismo y el usuario, al manejar, envía los datos de paso; los robots todavía no han llegado a ese punto. Si un robot teleoperado no genera un ROI suficiente en escenarios como aeropuertos, hoteles, restaurantes o comercio minorista, no puede formar el ciclo positivo de «el uso pagado genera datos, los datos mejoran la capacidad y la capacidad trae más uso». El ponente opina que Tesla quizá tenga los escenarios internos de fabricante automotriz y la capacidad organizacional para recorrer parte de esa vía, pero la mayoría de las empresas de robótica no puede simplemente copiarla.

\begin{knowledgebox}{Términos clave: vía real, teleoperación y vía sintética}
\begin{tabularx}{\textwidth}{p{0.20\textwidth}p{0.35\textwidth}X}
\toprule
Vía & Fuente de datos & Restricción principal \\
\midrule
Vía de flota real & Equipos producidos en serie que funcionan de manera natural y envían datos & Requiere que el producto ya sea útil y que la escala de hardware ya exista. \\
Vía de teleoperación & Una persona controla el robot a distancia y se registran las trayectorias de acción & Presión muy alta de personal, escenarios, regulación y ROI. \\
Vía de datos sintéticos & Generación masiva de tareas e interacciones en escenarios de simulación & Requiere Real2Sim de alta calidad, motores de física y sistemas de evaluación. \\
Vía híbrida & Poca calibración con datos reales y ampliación a gran escala con datos sintéticos & La dirección de ingeniería con más probabilidades de volverse dominante actualmente. \\
\bottomrule
\end{tabularx}
\end{knowledgebox}

\subsection{Resumen de la sección}

El cuello de botella de datos de la conducción autónoma se concentra principalmente en la cola larga y en los escenarios de seguridad; el de la inteligencia corporizada llega hasta la propia interacción física. Los robots no cuentan con un ciclo virtuoso de datos reales a gran escala ya establecido, por lo que los datos sintéticos y la simulación física no son un parche, sino el núcleo de la vía inicial.
\section{Physical Real2Sim: el umbral real de una buena simulación}

Tras explicar por qué la inteligencia corporizada debe apoyarse en la simulación, esta sección entra en la cuestión técnica más difícil: qué cuenta como una buena simulación. Xie Chen (谢晨) distingue entre el Real2Sim visual y el Real2Sim físico. El Real2Sim visual reconstruye la apariencia, la geometría y la escena tal como se ven; el Real2Sim físico, en cambio, debe trasladar a la simulación las propiedades mecánicas, las relaciones de interacción, el cuerpo del robot y los parámetros físicos del mundo real. Esta segunda versión es la que la inteligencia corporizada realmente necesita.

Una buena simulación requiere al menos cuatro componentes: activos y escenas sim-ready de alta calidad, un solver físico preciso, renderizado y simulación de sensores, y API, metadata y una cadena de herramientas que los algoritmos de entrenamiento puedan invocar. Si falta cualquiera de ellos, surgen problemas. Por ejemplo, si la puerta de un refrigerador se ve muy realista pero no se puede abrir, o si se abre pero con una fuerza incorrecta, la política que el robot aprenda en la simulación podría no transferirse.

\begin{figure}[H]
\centering
\includegraphics[width=0.96\textwidth]{real2sim-workflow.png}
\caption{Flujo de trabajo de Real2Sim: de la captura del mundo real a la reconstrucción, la calibración y una simulación reproducible. Diagrama conceptual propio, elaborado a partir de la conversación entre 00:32:41 y 00:46:18.}
\end{figure}

\begin{knowledgebox}{Lectura de la figura: Real2Sim convierte el mundo físico en un mundo entrenable}
Lo más importante de la figura es la «calibración», no la «reconstrucción». La captura real recupera la apariencia, la geometría, los parámetros mecánicos y las trayectorias de interacción; los activos de simulación deben poder ser tocados, jalados, empujados y sujetados por el robot; por último, hay que verificar mediante la evaluación en tareas reales si la política se transfiere. Solo al recorrer la cadena completa, Real2Sim deja de ser modelado 3D y se convierte en infraestructura de entrenamiento.
\end{knowledgebox}

\subsection{Cuatro criterios de una buena simulación}

Esta subsección concreta qué es una «buena simulación». Primero, los activos y las escenas deben contener suficiente información física, incluidos cuerpos de colisión, articulaciones, fricción, peso, par y materiales. Segundo, el solver subyacente debe ser lo bastante preciso y capaz de resolver problemas distintos, como cuerpos rígidos, cuerpos no rígidos, telas y líquidos. Tercero, la salida de la simulación debe servir al entrenamiento de algoritmos mediante API y metadata, no solo mostrarse a personas. Cuarto, la simulación debe ser lo bastante eficiente para dar servicio al aprendizaje por refuerzo en paralelo, sobre todo al RL con visión en el lazo.

Hay aquí un punto que se pasa por alto con facilidad: el cuerpo del robot en la simulación también tiene que ser preciso. En la entrevista se menciona que la mano del robot real de un cliente podía cargar uno o dos kilogramos, pero en la simulación solo podía cargar 0.1 kilogramos; el equipo pasó seis meses sin lograr ajustarlo. Esto muestra que la Sim2Real gap no proviene solo de la escena, sino también del propio modelo del robot. Una mano errónea, un límite articular erróneo o una curva de par errónea bastan para distorsionar los resultados del entrenamiento.

\begin{importantbox}{Criterios de una buena simulación}
Una buena simulación no es la que tiene buena imagen, sino la que permite llevar a robots reales los algoritmos entrenados en ella. Debe contar con parámetros físicos reales, un solver preciso, una cadena de herramientas entrenable, alta capacidad de paralelización y un ciclo de evaluación que verifique continuamente el éxito o el fracaso de Sim2Real.
\end{importantbox}

\begin{knowledgebox}{Términos clave: Real2Sim físico}
\begin{tabularx}{\textwidth}{p{0.20\textwidth}p{0.35\textwidth}X}
\toprule
Término & Significado & Por qué afecta el entrenamiento \\
\midrule
Solver & Resolvedor del motor físico que calcula el movimiento, las colisiones y las restricciones & Determina si los resultados mecánicos son confiables. \\
Collision Body & Geometría que participa en el cálculo de colisiones & Sin cuerpos de colisión correctos, sujetar objetos y abrir puertas deja de ser realista. \\
Sim-ready Asset & Activo que puede usarse directamente en interacciones de simulación & No basta con la apariencia del modelo; también debe incluir parámetros físicos e interfaces. \\
Metadata & Etiquetas, estados, parámetros e información semántica adicionales & El entrenamiento y la evaluación necesitan leer esta información estructurada. \\
Parallel Environment & Múltiples instancias de simulación que se ejecutan en paralelo en RL & Determina el rendimiento del entrenamiento y la eficiencia de exploración. \\
\bottomrule
\end{tabularx}
\end{knowledgebox}

\subsection{Diferencias entre el aprendizaje por imitación y el aprendizaje por refuerzo}

Esta subsección examina los requisitos de la simulación desde el objetivo de entrenamiento. El aprendizaje por imitación puede usar datos de lazo abierto: una persona o un sistema de teleoperación demuestra una trayectoria y el modelo aprende a replicarla. Es adecuado para demos tempranas, pero su generalización suele ser limitada. Por ejemplo, el modelo puede aprender a abrir la tapa de cierta botella alta, pero no necesariamente generaliza a botellas bajas, a otras formas de girar o a otros materiales. El aprendizaje por refuerzo, en cambio, necesita un entorno de lazo cerrado en el que el robot pruebe y se equivoque una y otra vez y aprenda de la reward, lo que exige que la simulación sea a la vez precisa y eficiente.

El ponente compara el RL de robots con el posentrenamiento de los modelos grandes: primero se tiene un modelo base y luego se hace fine-tune mediante aprendizaje por refuerzo. La inteligencia corporizada también podría necesitar un VLA o un modelo base robótico, cuya política se ajuste después mediante RL en una simulación de alta calidad. La dificultad es que los entornos de RL con visión en el lazo son muy pesados: cuántos environment pueden ejecutarse en paralelo en una GPU y qué tan alto es el total FPS determinan directamente si el entrenamiento puede escalar.

\begin{warningbox}{Si la simulación no admite RL, difícilmente sostendrá la generalización de la siguiente etapa}
Una simulación que solo sirve al aprendizaje por imitación puede generar trayectorias y demostraciones; una simulación que sirve al RL debe permitir que la política actúe, observe, falle, reintente y obtenga reward. Esta última exige mucho más en precisión física, eficiencia en paralelo y cadena de herramientas.
\end{warningbox}

\subsection{Simulación con calidad de agua potable}

Esta subsección resume la visión de Guanglun Zhineng (光轮智能): una buena simulación debería ser tan accesible como el agua potable. La metáfora es contundente, porque convierte la simulación de «herramienta misteriosa de algún laboratorio» en infraestructura pública de la industria. Hoy muchas simulaciones son como agua sin filtrar: parece haber mucha, pero beberla causa problemas; lo que de verdad impulsará a la industria en el futuro es que todo investigador en robótica pueda acceder a entornos de simulación de calidad suficiente, suficientemente utilizables y suficientemente verificables.

Para lograrlo, no basta con construir un simulador. En la ronda final de preguntas rápidas, Xie Chen subraya que «simulación no equivale a simulador»: el simulador es solo la capa del motor físico, mientras que la simulación también abarca activos, escenas, renderizado, API, framework, metadata, control de calidad y evaluación. Muchas empresas empiezan directamente por el simulador de bajo nivel y, en consecuencia, pueden pasar por alto los activos y las escenas, que son lo que realmente escasea. El camino de Guanglun es construir primero activos, escenas y cadenas de herramientas de alta calidad, y solo al final profundizar en el simulador de bajo nivel.

\begin{importantbox}{Simulación no equivale a simulador}
El simulador es una condición necesaria pero no suficiente. Lo que realmente le falta a la inteligencia corporizada son activos físicos interactivos, escenas escalables, interfaces entrenables, estándares de calidad verificables y un ciclo cerrado capaz de llevar de vuelta al sistema los fallos reales.
\end{importantbox}

\subsection{Resumen de la sección}

Physical Real2Sim es el núcleo técnico de EP109. Exige que el equipo lleve a la simulación la información física del mundo real con bajo costo y a escala, y que la calibre continuamente con los resultados del entrenamiento de modelos y del despliegue real. El criterio de una buena simulación no es el realismo visual, sino si la política entrenada funciona en el mundo real.
\section{Empresas de datos, Meta/Scale y organización del talento}

Lo anterior trató de la tecnología de datos para robots; esta sección se ocupa de por qué la capacidad de producir datos se ha vuelto una capacidad estratégica de las grandes empresas. En la entrevista, Zhang Xiaojun (张小珺) y Xie Chen (谢晨) comentan la adquisición de Scale AI por parte de Meta a un precio elevado y el papel de Alexandr Wang. Lo importante aquí no es el chisme de la operación en sí, sino una conclusión más amplia: en la siguiente etapa, la competencia en IA no dependerá solo del poder de cómputo y de los modelos, sino también de la capacidad de organizar la producción de datos de alta calidad.

En la entrevista se explica que el valor de Scale consiste en que «los datos se venden cada vez más caros», porque las personas que emplea son cada vez mejores, las tareas son cada vez más complejas y los clientes necesitan cada vez más demostraciones y control de calidad de alto nivel. Para los modelos grandes y para los robots, los buenos datos ya no son solo etiquetado de bajo costo, sino el resultado conjunto de talento de primer nivel, procesos, relationship, sistemas de inspección de calidad y conocimiento del cliente. Si Meta quiere controlar la materia prima básica de la competencia futura en IA, buscará incorporar esta capacidad de datos a su propia estructura mediante una compra.

\begin{figure}[H]
\centering
\includegraphics[width=0.94\textwidth]{meta-scale-data.png}
\caption{La señal de la adquisición de Scale AI por parte de Meta: las empresas de diez billones del futuro necesitan dominar la capacidad de datos para IA. Diagrama conceptual propio, elaborado a partir de la conversación de 00:48:55--00:53:57.}
\end{figure}

\begin{knowledgebox}{Lectura de la figura: por qué la capacidad de datos pasó de ser subcontratada a ser estratégica}
En la figura se distinguen tres niveles de valor: el proceso de producción de datos, la demostración humana y la red de talento. Al principio, el etiquetado podía subcontratarse; cuando las tareas pasan a ser agentic, de robótica, de razonamiento complejo y de entornos de RL de alta calidad, la producción de datos se parece cada vez más a una capacidad central de investigación y desarrollo. La señal de Meta/Scale es esta: quien logre organizar a las mejores personas para definir los datos, filtrarlos y elevar su nivel de exigencia (bar) estará más cerca de la capacidad de los modelos de la siguiente etapa.
\end{knowledgebox}

\subsection{Los buenos datos son el producto de la diversidad de escenarios y de demostraciones de alta calidad}

Esta sección parte de una afirmación central: los buenos datos son el producto de dos cosas; la primera, escenarios diversos; la segunda, demostraciones humanas de alta calidad. Para los modelos grandes, la diversidad puede consistir en escenarios de tareas como matemáticas, programación, derecho o medicina; la demostración humana puede ser los pasos para resolver un problema, los criterios de evaluación o preguntas de examen de alta calidad. Para los robots, la diversidad está en hogares, hoteles, fábricas, restaurantes, anaqueles, cocinas, manijas de puertas y distintos objetos; la demostración humana, en cambio, consiste en trayectorias de movimiento, diseño de tareas, juicio sobre los fallos y diseño de la reward.

Además, el ponente sostiene que el mejor maestro no solo demuestra, sino que plantea problemas. Esta analogía explica por qué el trabajo con datos es cada vez más caro: no se trata de pedirle a alguien que haga una demostración cualquiera, sino de que diseñe problemas capaces de inducir el progreso del modelo, proporcione feedback de alta calidad y convierta los fallos en tareas de entrenamiento para la siguiente ronda. En el caso de los robots, un buen equipo de datos debe entender los escenarios físicos y también el entrenamiento de modelos, y además debe ser capaz de operar procesos de colaboración entre personas y máquinas a gran escala.

\begin{importantbox}{Fórmula de la calidad de los datos}
En el contexto de este episodio, los datos de alta calidad pueden entenderse así:
\[
\mathrm{Data\ Quality}\approx \mathrm{Diversity\ of\ Scenarios}\times \mathrm{Quality\ of\ Human\ Guidance}
\]
La diversidad de escenarios determina cuánto del mundo ha visto el modelo, y la guía humana de alta calidad determina si esos mundos pueden convertirse en señales de aprendizaje efectivas. Si falta cualquiera de las dos, los datos se empobrecen.
\end{importantbox}

\subsection{Cuellos de botella actuales de los datos sintéticos}

Esta sección vuelve a los datos sintéticos propiamente dichos. Xie Chen reconoce que, desde los primeros principios, el Sim2Real gap siempre existirá, porque lo sintético y lo real son distintos por naturaleza. Lo que importa no es eliminar el gap, sino saber a partir de qué magnitud del gap los datos resultan útiles y cómo reducirlo de manera continua. Esto exige que una empresa de datos no se limite a entregar datos, sino que también pueda llevar los algoritmos a la práctica, hacer despliegues reales y cerrar el ciclo de evaluación; de lo contrario, nunca sabrá dónde están las carencias.

Los cuellos de botella de los datos sintéticos pueden dividirse en tres etapas. En la etapa Real2Sim hay que aumentar la diversidad y la precisión de la información obtenida del mundo real; en la etapa Simulation hay que mejorar el motor físico subyacente, los recursos, los escenarios y la eficiencia en paralelo; en la etapa Sim2Real hay que mejorar la generalización y la efectividad de las políticas entrenadas. Por eso Guanglun (光轮), al mismo tiempo que produce datos, necesita entender los algoritmos de robótica de sus clientes y su puesta en práctica real.

\begin{warningbox}{Una empresa de datos no puede limitarse a entregar archivos}
Si una empresa de datos sintéticos solo se encarga de generar datos y el cliente entrena el modelo por su cuenta, a la empresa de datos le resultará difícil saber dónde está el gap real. Para mejorar de forma continua la calidad de los datos, debe participar en la evaluación del modelo, el despliegue real y el análisis de fallos.
\end{warningbox}

\subsection{El dilema de Physical Intelligence}

Esta sección examina la valoración que se hace en la entrevista de Physical Intelligence. Xie Chen considera que PI va muy adelante en modelos fundacionales corporizados, pero enfrenta un dilema respecto a la simulación: su discurso externo enfatiza los datos reales, mientras que internamente es muy consciente de que la simulación es necesaria. Esta contradicción afecta la atracción de talento, porque a los mejores especialistas en simulación les cuesta unirse a una organización que públicamente no da importancia a la simulación.

El valor didáctico de esta discusión es el siguiente: la simulación no es un pequeño componente de investigación, sino la combinación de cuatro cosas: tecnología, ingeniería, algoritmos y operación. Los equipos de investigación suelen dominar los algoritmos y los artículos, pero la simulación a gran escala requiere además plataformas de ingeniería, procesos operativos, inspección de calidad humana y captura física. Precisamente porque es pesada, sucia y compleja, se convierte en una oportunidad independiente dentro de la cadena industrial.

\begin{knowledgebox}{Términos clave: las cuatro cosas de la capacidad de simulación}
\begin{tabularx}{\textwidth}{p{0.18\textwidth}p{0.36\textwidth}X}
\toprule
Capacidad & Contenido concreto & Por qué la investigación por sí sola no basta \\
\midrule
Tecnología & Motor físico, renderizado, recursos y escenarios & Requiere acumulación de base a largo plazo. \\
Ingeniería & Plataformas a gran escala, cadena de herramientas, API y eficiencia en paralelo & Tanto el entrenamiento como las entregas al cliente exigen sistemas estables. \\
Algoritmos & Real2Sim, Sim2Real, RL y evaluación de calidad & Hay que convertir los datos en beneficios para el modelo. \\
Operación & Captura con humanos en el ciclo, inspección de calidad, retroalimentación de clientes y estandarización & La información del mundo físico no puede obtenerse de forma totalmente automática. \\
\bottomrule
\end{tabularx}
\end{knowledgebox}

\subsection{Resumen de la sección}

La discusión sobre Meta/Scale extiende el EP109 de los robots a toda la industria de la IA: la competencia futura girará cada vez más en torno a la capacidad de producir datos. Esto vale en especial para la inteligencia corporizada, porque los datos de alta calidad requieren diversidad del mundo y, además, demostraciones humanas, parámetros físicos, plataformas de ingeniería y un ciclo cerrado de evaluación.
\section{Mapping de la cadena industrial global de la inteligencia corporizada}

Esta sección organiza el mapa de la cadena industrial que presentó Xie Chen (谢晨). La inteligencia corporizada se divide en varios tipos de empresas porque debe inventar al mismo tiempo «la plataforma del vehículo de próxima generación» y «los algoritmos L4/L5 que corren sobre esa plataforma». La conducción autónoma cuenta al menos con una plataforma madura, el automóvil; en la inteligencia corporizada, en cambio, la plataforma, el cuerpo robótico, los sensores, las manos, los modelos, los datos y los escenarios están todos en una etapa temprana. Cuando la dificultad es tan alta, la cadena industrial tiende a dividir el trabajo.

En el programa se distinguen, en términos generales, cuatro tipos de empresas: empresas de hardware, empresas de Foundation Model, empresas que integran software y hardware para escenarios verticales, y empresas de infraestructura centradas en la simulación. Esta clasificación no es una etiqueta estática, sino que indica dónde coloca el riesgo cada empresa: las de hardware resuelven el cuerpo robótico y la cadena de suministro; las de modelos, el cerebro y la recipe de datos; las de implantación, el ROI del escenario; y las de simulación, los datos y el mundo de entrenamiento.

\begin{figure}[H]
\centering
\includegraphics[width=0.96\textwidth]{embodied-supply-chain-map.png}
\caption{Cadena industrial global de la inteligencia corporizada: división del trabajo entre empresas de hardware, de modelos base, de implantación vertical y de simulación. Diagrama conceptual propio, elaborado a partir de la conversación de 00:55:25--01:09:22.}
\end{figure}

\begin{knowledgebox}{Lectura de la figura: cuatro tipos de empresas corresponden a cuatro recursos escasos}
En la capa de hardware escasean un cuerpo robótico confiable, el costo y la cadena de suministro; en la capa de modelos escasean el cómputo, el talento y una recipe de datos de propósito general; en la capa de implantación vertical escasean los clientes, los escenarios y la operación; en la capa de simulación escasean un mundo físico de alta calidad, la cadena de herramientas y el ciclo cerrado de evaluación. La inteligencia corporizada es demasiado pesada, por lo que a corto plazo se parece más a un ecosistema industrial que a una sola empresa que lo abarque todo.
\end{knowledgebox}

\subsection{Hardware, modelos, implantación y simulación}

Esta subsección desarrolla uno por uno los cuatro tipos de empresas. Entre las empresas de hardware destacan Yushu (宇树) y las empresas de manos diestras, que llevaron el cuerpo robótico a la investigación y a la industria y, en cierta medida, establecieron un estándar de investigación. Las empresas de Foundation Model incluyen Physical Intelligence, Skild, NVIDIA/GEAR, Google/DeepMind, entre otras, que buscan el «cerebro» del robot y la scaling recipe. Las empresas de implantación que integran software y hardware incluyen Figure, Tesla Optimus, The Bot Company, DynaRobotics, entre otras, y se enfocan en escenarios verticales concretos. Las empresas de simulación, como Guanglun (光轮) y Genesis, construyen infraestructura en torno a Real2Sim, Simulation y Sim2Real.

La valoración que hace el ponente de Yushu resulta muy ilustrativa para la industria: ofrece el cuerpo robótico a la comunidad académica internacional, de modo que los artículos y los estudiantes adquieren el hábito de trabajar con robots de Yushu, y luego trasladan ese hábito a la industria. Esto muestra que el «estándar» en la capa de hardware no proviene solo de las especificaciones, sino también del ecosistema de desarrolladores, del ecosistema de publicaciones y de la movilidad del talento. En la capa de simulación, Guanglun aspira a ser ese mismo tipo de estándar: no un proveedor más, sino la base común de la simulación corporizada.

\begin{knowledgebox}{Tabla de roles de la cadena industrial}
\begin{tabularx}{\textwidth}{p{0.18\textwidth}p{0.32\textwidth}X}
\toprule
Rol & Tarea central & Principales dimensiones de competencia \\
\midrule
Empresas de hardware & Cuerpo robótico, manos, sensores y equipo de captura & Costo, confiabilidad, ecosistema de desarrolladores, cadena de suministro. \\
Empresas de modelos & Modelos base para robots, cerebro, VLA y scaling recipe & Cómputo, talento, pirámide de datos, capacidad de RL. \\
Empresas de implantación & Entrega para escenarios como fabricantes de automóviles, almacenes, hogares y restaurantes & ROI del cliente, integración de software y hardware, operación y seguridad. \\
Empresas de simulación & Generar mundos entrenables que sostengan el ciclo cerrado de datos y algoritmos & Activos físicos, solver, API, tasa de éxito de Sim2Real. \\
\bottomrule
\end{tabularx}
\end{knowledgebox}

\subsection{Diferencias de oportunidad entre China y Estados Unidos}

La subsección anterior trató la cadena industrial global; esta examina la estructura por país. Xie Chen considera que en Estados Unidos hay oportunidades de emprendimiento en la capa de modelos, porque su ecosistema pone más énfasis en la división industrial del trabajo y los usuarios finales tienen mayor disposición a pagar por software, inferencia y servicios de automatización. Empresas como OpenAI y NVIDIA pueden obtener ingresos muy altos mediante modelos, cómputo y ecosistema, por lo que la capa de modelos puede constituir empresas independientes.

La estructura de China es distinta. En el país, la disposición a pagar por software y por tokens es relativamente baja, mientras que las capacidades de hardware, cadena de suministro y equipo completo son más fuertes, de modo que muchas empresas siguen la ruta de integrar software y hardware, entregar soluciones de extremo a extremo y vender hardware. En el programa se menciona que empresas como Zijie (字节), Xiaomi (小米) y Lixiang (理想) podrían ser más adecuadas para desarrollar el «cerebro», porque cuentan con cómputo, talento en IA, puntos de acceso de producto, cadena de suministro o escenarios de vehículos o robots. Este juicio no significa que China carezca de capacidad en modelos, sino que el modelo de negocio y la organización industrial colocarán esa capacidad dentro de sistemas integrados más verticales.

\begin{figure}[H]
\centering
\includegraphics[width=0.94\textwidth]{china-us-embodied-opportunity.png}
\caption{Diferencias de oportunidad en inteligencia corporizada entre China y Estados Unidos: Estados Unidos tiene más oportunidades en la capa de modelos; China es más adecuada para integrar cerebro + cuerpo robótico + escenario. Diagrama conceptual propio, elaborado a partir de la conversación de 01:09:22--01:15:33.}
\end{figure}

\begin{knowledgebox}{Lectura de la figura: las diferencias provienen de la estructura de pago y de la división industrial del trabajo}
A la izquierda, en Estados Unidos es más fácil que surja una capa de modelos independiente, porque el pago por software, el pago por inferencia y la división del trabajo en el ecosistema permiten que las empresas de modelos obtengan ingresos. A la derecha, en China es más fácil integrar cerebro, cuerpo robótico y escenario, porque la cadena de suministro de hardware es fuerte, el cobro por software por separado es débil y las empresas de producto necesitan más un valor de extremo a extremo. Al leer la figura, no debe entenderse como una diferencia de nivel tecnológico, sino como una diferencia de estructura comercial.
\end{knowledgebox}

\subsection{La pirámide de datos y la convergencia de caminos distintos}

Esta subsección explica por qué rutas técnicas distintas podrían terminar convergiendo. En apariencia, algunos equipos dicen seguir la ruta de datos reales, otros la de simulación, unos ponen el énfasis en imitation learning y otros en RL. Sin embargo, Xie Chen observa que los clientes de primer nivel adoptan cada vez más, internamente, una data pyramid: en la base, datos de internet o de propósito general; en el medio, datos sintéticos y simulación; en la cima, una cantidad pequeña de datos reales de alta calidad.

Esta estructura piramidal concuerda con la lógica expuesta antes. Los datos de internet aportan conocimiento del mundo y preentrenamiento visual y de lenguaje; los datos sintéticos aportan entornos de tareas controlables, diversos y paralelizables; los datos reales aportan grounding, calibración y la validación final. Los equipos todavía difieren en RL, VLA, proporción de datos sintéticos y arquitectura de algoritmos, pero el «co-training real + sintético» ya es un consenso entre las instituciones principales.

\begin{importantbox}{Pirámide de datos}
Es probable que la inteligencia corporizada no dependa solo de datos reales ni solo de datos de simulación. La ruta más sólida es: en la base, datos de propósito general que construyen conocimiento del mundo; en el nivel medio, datos sintéticos o de simulación que amplían la distribución de tareas; en la cima, datos reales para calibración, demostración y validación del despliegue.
\end{importantbox}

\subsection{Resumen de la sección}

La cadena industrial de la inteligencia corporizada mantendrá la división del trabajo por mucho tiempo. El hardware, los modelos, la implantación y la simulación controlan cada uno recursos escasos distintos; las diferencias entre China y Estados Unidos provienen más del modelo de negocio y de la organización industrial; y las distintas rutas podrían converger finalmente en la pirámide de datos y en el entrenamiento mixto con datos reales y sintéticos.
\section{NVIDIA is a simulation company}

La sección anterior trató la división del trabajo en la cadena industrial; esta se centra en NVIDIA. Xie Chen (谢晨) recuerda que Jensen Huang (黄仁勋) dijo internamente: NVIDIA is a simulation company. El sentido de esta frase no es solo publicitario: replantea a NVIDIA, de una empresa que «vende GPU», a una empresa que «usa el cómputo para crear mundos en los que se puede entrenar». Los videojuegos son simulación al servicio de la experiencia humana; Omniverse es simulación al servicio de la industria y de los gemelos digitales; Isaac Sim es simulación al servicio de los robots y de la IA física.

El three computer problem de NVIDIA también se relaciona con esto: la primera computadora es el centro de datos; la segunda, el dispositivo de IA física en el extremo; la tercera, la computadora de simulación. El centro de datos entrena los grandes modelos; los robots y vehículos en el extremo ejecutan los modelos; la computadora de simulación genera el mundo físico, los datos y los entornos de evaluación. Cuanto más grave es la escasez de datos de robótica, más se parece la tercera computadora a una puerta de entrada estratégica.

\begin{figure}[H]
\centering
\includegraphics[width=0.96\textwidth]{nvidia-simulation-company.png}
\caption{NVIDIA is a simulation company: la capacidad de cómputo, los motores de física, los entornos de simulación y el entrenamiento de robots forman un ciclo cerrado. Diagrama conceptual propio, elaborado a partir de la conversación de 01:15:33--01:21:25.}
\end{figure}

\begin{knowledgebox}{Lectura de la figura: la estrategia de simulación de NVIDIA conecta tres tipos de cómputo}
A la izquierda de la figura está la capacidad de cómputo para entrenamiento e inferencia; en el centro, los motores de física, el renderizado, los activos y las API; a la derecha, el despliegue en robots y en conducción autónoma. La fortaleza de NVIDIA no es una herramienta aislada, sino unir en una infraestructura las GPU, Omniverse, Isaac Sim, PhysX, Warp/Newton y el ecosistema de robótica. Cuanto más se convierte la simulación en la forma de producir datos, más se acerca una empresa de cómputo a ser una empresa de datos.
\end{knowledgebox}

\subsection{Isaac Sim, MuJoCo y la siguiente generación de motores de física}

Esta subsección estructura un poco el ecosistema de simuladores que aparece en la entrevista. Isaac Sim, de NVIDIA, está construido sobre PhysX y el renderizado de Omniverse; sus ventajas son las GPU, el renderizado y el ecosistema industrial. MuJoCo, vinculado a Google/DeepMind, ha sido durante mucho tiempo un motor de física de uso común en la investigación en robótica; sus ventajas son la física y la API, pero el renderizado no es su punto fuerte. Como el RL de extremo a extremo requiere un paralelismo masivo en GPU, la versión de MuJoCo para CPU difícilmente alcanza el caudal de procesamiento necesario; MJX y MuJoCo Warp/Newton representan la evolución hacia la aceleración por GPU, el ecosistema de código abierto y una base física unificada.

Lo que más vale la pena aprender aquí no es qué motor gana, sino que el simulador se convertirá en una capa de la competencia entre ecosistemas. Cuanto más abierto, eficiente y mantenible sea el motor de física subyacente, más rápido podrán desarrollarse las capas superiores: activos, API, generación de datos y plataformas de RL. Si una empresa como Guanglun (光轮) logra aportar activos, parámetros físicos y cadenas de herramientas, se volverá una parte importante de este ecosistema sin necesidad de poseer en exclusiva el motor subyacente.

\begin{knowledgebox}{Términos clave: las cuatro capas de la pila de simulación}
\begin{tabularx}{\textwidth}{p{0.18\textwidth}p{0.36\textwidth}X}
\toprule
Capa & Contenido representativo & Función \\
\midrule
Physics Solver & PhysX, MuJoCo, Newton, etc. & Calcula la mecánica, las colisiones, las restricciones y el movimiento. \\
Rendering & Cadenas de renderizado como Omniverse & Genera las observaciones visuales y las entradas de los sensores. \\
Sim-ready Assets & Activos interactivos, escenas, cuerpos de robot & Proporcionan la materia prima del mundo de entrenamiento. \\
API/Framework & Exportación de datos, metadata, interfaces de entornos de RL & Permiten que el entrenamiento y la evaluación de modelos invoquen la simulación. \\
\bottomrule
\end{tabularx}
\end{knowledgebox}

\subsection{Datos sintéticos y grandes modelos}

Esta subsección pasa de los robots a los grandes modelos. En la entrevista, Xie Chen sostiene que los grandes modelos de lenguaje natural también dependen cada vez más de datos sintéticos, porque los datos de internet tienden a agotarse y el propio modelo puede generar nuevas tareas, nuevas soluciones y nuevas evaluaciones. En cierto sentido, GPT también es un generador de datos sintéticos: produce lenguaje natural sin cesar, que puede usarse para entrenar, destilar o evaluar otros modelos.

Para los grandes modelos corporeizados, en cambio, los datos sintéticos son más complejos: no solo requieren texto, sino también 3D, la perspectiva del robot, múltiples sensores, trayectorias de acción e interacción física. Los datos sintéticos de lenguaje natural pueden apoyarse en la generación y la verificación por modelos; los datos sintéticos corporeizados tienen que apoyarse en mundos simulados y en la calibración con la física real. Por eso EP109 da a los datos sintéticos más peso que la discusión habitual sobre grandes modelos.

\begin{warningbox}{Los datos sintéticos de lenguaje y los corporeizados no tienen la misma dificultad}
Los datos de lenguaje pueden generarse y revisarse principalmente en el espacio simbólico; los datos corporeizados tienen que pasar por las restricciones del mundo físico. Para saber si la trayectoria de un robot es válida no basta con la lógica textual: hay que considerar el contacto, la mecánica, la seguridad, los actuadores y el despliegue real.
\end{warningbox}

\subsection{Resumen de la sección}

La estrategia de simulación de NVIDIA muestra que la infraestructura de IA del futuro no se limita a los clústeres de entrenamiento y a los chips en el extremo, sino que también incluye la computadora de simulación que genera los mundos de entrenamiento. Para la robótica, la simulación es a la vez una herramienta de producción de datos, un entorno de evaluación de modelos y una puerta de entrada al ecosistema industrial.
\section{El escenario final: entre universos, entre mundos, entre cuerpos}

Esta sección analiza el modelo del escenario final, más lejano. Xie Chen (谢晨) considera que el modelo final debe ser cross universe, cross world, cross embodiment, es decir, capaz de transferirse entre universos, entre mundos y entre cuerpos. La formulación suena grandiosa, pero detrás hay un objetivo sencillo: mejorar la generalización. Los seres humanos pueden pasar con rapidez de la realidad a un juego, de la cocina a la oficina, de la mano a una herramienta; si un robot solo puede trabajar en una simulación, en una habitación o con un cuerpo, todavía no es inteligencia general.

Los datos de videojuegos tienen aquí un significado especial. Los juegos no son necesariamente fieles a la física, pero ofrecen mundos ricos, reglas, objetivos e interacción en primera persona, por lo que resultan adecuados para mejorar, en la etapa de preentrenamiento, la capacidad del agent de transferirse entre mundos. DeepMind y otras instituciones usan grandes cantidades de datos de juegos precisamente porque el entrenamiento entre universos acostumbra al modelo a actuar con reglas y distribuciones de observación distintas. No sustituye el Real2Sim físico, pero complementa una world diversity más amplia.

\begin{figure}[H]
\centering
\includegraphics[width=0.92\textwidth]{cross-universe-model.png}
\caption{Entre universos, entre mundos, entre cuerpos: el modelo del escenario final debe mejorar la generalización entre entornos. Diagrama conceptual propio, elaborado a partir de la conversación de 01:21:25--01:23:28.}
\end{figure}

\begin{knowledgebox}{Lectura de la figura: cada tipo de transferencia resuelve un problema de generalización distinto}
La transferencia entre universos resuelve el domain shift entre simulación, juegos y realidad; la transferencia entre mundos resuelve las diferencias entre entornos como hogares, fábricas, hoteles y restaurantes; la transferencia entre cuerpos resuelve las diferencias entre los cuerpos de los robots y sus espacios de acción. Al leer la figura, hay que notar que las tres no son consignas, sino tres tipos de diferencias en la distribución de los datos. El modelo del escenario final debe ser capaz de extraer la esencia de la tarea a través de estas distribuciones.
\end{knowledgebox}

\begin{knowledgebox}{Términos clave: los tres tipos de transferencia}
\begin{tabularx}{\textwidth}{p{0.18\textwidth}p{0.36\textwidth}X}
\toprule
Capacidad & Significado & Requisitos de entrenamiento \\
\midrule
Entre universos & Transferirse entre juegos, simulación y el mundo real & Datos de mundos diversos, aprendizaje de reglas y domain adaptation. \\
Entre mundos & Transferirse entre distintos entornos de tarea & Los recursos de escena, la distribución de tareas y la retroalimentación real deben ser suficientemente amplios. \\
Entre cuerpos & Transferirse entre distintos cuerpos de robot & Hay que modelar el espacio de acción, los sensores, la geometría del cuerpo y las restricciones de control. \\
\bottomrule
\end{tabularx}
\end{knowledgebox}

\subsection{Todavía en la etapa GPT-1}

Esta subsección vuelve al diagnóstico de la etapa actual. Xie Chen considera que la inteligencia corporizada en su conjunto sigue en la etapa GPT-1; esto no significa que no tenga ninguna capacidad, sino que aún no se ha encontrado una receta estable de scaling law. La analogía con el FSD de Tesla resulta útil: al principio, añadir datos de forma continua no siempre producía mejoras sostenidas, hasta que se consolidó la vía de extremo a extremo; solo entonces la expansión de datos, cómputo y modelo empezó a traducirse de manera más estable en mejoras de capacidad. La inteligencia corporizada todavía no ha vivido ese momento claro.

Sin embargo, no es pesimista, por tres razones: la densidad de equipos fundadores y de científicos que hoy entran en la inteligencia corporizada es mucho mayor que en los inicios de la conducción autónoma; la atención del capital y de la industria es mayor; y los modelos grandes y la conducción autónoma ya aportaron experiencia en scaling, transformer, volante de datos, RL y entrenamiento de extremo a extremo. Dicho de otro modo, es una etapa temprana, pero no una etapa temprana que empiece desde cero.

\begin{figure}[H]
\centering
\includegraphics[width=0.96\textwidth]{embodied-gpt1-stage.png}
\caption{La inteligencia corporizada sigue en la etapa GPT-1: aún no se ha encontrado una scaling law estable y la recipe de datos sigue en exploración. Diagrama conceptual propio, elaborado a partir de la conversación de 01:23:28--01:28:21.}
\end{figure}

\begin{knowledgebox}{Lectura de la figura: la etapa GPT-1 significa que la recipe no ha convergido}
Lo más importante de la figura es «unknown recipe». La industria sabe que necesita datos, cómputo, modelos, simulación y RL, pero todavía no sabe qué combinación mejorará de forma estable la capacidad de los robots reales. Más que indicar que la capacidad es débil, la etapa GPT-1 indica que aún no se ha formado una vía de escalamiento tan clara como la de GPT-3/ChatGPT.
\end{knowledgebox}

\begin{importantbox}{Scaling law moment}
El scaling law moment de la inteligencia corporizada no es una ronda de financiamiento, una demo ni el lanzamiento de un modelo aislado, sino el momento en que el equipo descubre que, con una arquitectura y una canalización de datos dadas, seguir añadiendo datos de alta calidad, cómputo y entornos mejora de forma predecible la capacidad en tareas reales.
\end{importantbox}

\subsection{El progress es más sano que el financiamiento}

Esta subsección conserva el juicio de emprendedor de la segunda mitad de la entrevista. A Xie Chen le preocupa que la industria se fije demasiado en los montos de financiamiento y en las valuaciones, y descuide el progress real. Desde la perspectiva de RL, si el reward model de un equipo se convierte en «dar conferencias, hacer demos y atraer inversionistas», su conducta queda guiada por una recompensa equivocada; un reward model más sano consiste en atender a los clientes, crear valor, cobrar y mejorar el producto de forma continua con la retroalimentación de los clientes.

Este pasaje es, en realidad, coherente con la línea técnica principal: ya se trate de datos, de simulación o de una empresa, todo necesita retroalimentación real. Para el modelo, es la retroalimentación del despliegue real; para la empresa, la retroalimentación de los pagos y las recompras de los clientes; para la industria, el progress real y no el narrative progress. En la entrevista, el paso de Guanglun (光轮) de la conducción autónoma a la inteligencia corporizada también implicó pasar de «yo sé de simulación» a «volver a aprender la inteligencia corporizada desde el nivel de licenciatura», lo que muestra que la experiencia de problemas de baja dimensión no puede aplicarse directamente a problemas de alta dimensión.

\begin{warningbox}{Una señal de burbuja en la industria}
Si un equipo optimiza sobre todo la narrativa para el financiamiento, la impresión que causan las demos y la valuación, en lugar del valor para el cliente, la mejora del modelo y el despliegue real, su reward model está desviado. En una industria con tanta carga de ingeniería como la inteligencia corporizada, una recompensa equivocada es más peligrosa que un fracaso técnico de corto plazo.
\end{warningbox}

\subsection{Resumen de la sección}

El escenario final de la inteligencia corporizada es la generalización entre universos, entre mundos y entre cuerpos; en la etapa actual todavía se busca la scaling law. Quizá la industria encuentre pronto la recipe, pero a condición de alinear el reward con el progress real: mejoras reales del modelo, valor real para el cliente y retroalimentación real del despliegue.
\section{Síntesis y ampliación}

Esta sección resume toda la entrevista en un marco reutilizable. Lo esencial de EP109 no es que «las empresas de simulación son importantes», sino explicar por qué la IA del mundo físico necesita una nueva forma de producir datos. Los datos de internet impulsaron el despegue de los modelos de lenguaje y los datos de las flotas de vehículos formaron un ciclo que se refuerza a sí mismo en la conducción autónoma. La inteligencia corporizada no cuenta con datos del mundo ya disponibles, por lo que tiene que construir activamente su mundo de entrenamiento con Real2Sim, datos sintéticos, entornos de simulación, una pequeña cantidad de retroalimentación real y sistemas de evaluación.

Desde el punto de vista técnico, los buenos datos sintéticos deben pasar tres validaciones. Primera: ¿cubren las variables más escasas, más peligrosas y con mayor efecto sobre el modelo en las tareas reales? Segunda: ¿puede demostrarse mediante el sistema de evaluación que mejoran el modelo? Tercera: ¿pueden calibrarse durante el despliegue en robots reales? Los datos que solo pasan la primera prueba son una biblioteca de escenarios; los que pasan la segunda son datos de entrenamiento, y solo los que pasan la tercera se acercan a una infraestructura que puede escalar.

\begin{importantbox}{Ubicar EP109 en la serie de IA de Zhang Xiaojun (张小珺)}
EP109 forma una línea continua con EP121 sobre los robots de DeepMind, EP132 sobre los robots de Xinghaitu (星海图) y EP134, la revisión general sobre datos: EP121 trata los modelos de robots y la transferencia entre distintos cuerpos robóticos; EP132, el robot completo y la Data Recipe; EP134, el valor de los datos para la IA, y EP109 desglosa la infraestructura de producción de datos de la inteligencia corporizada en Real2Sim, simulación, datos sintéticos, evaluación y división del trabajo en la industria.
\end{importantbox}

\subsection{Cinco takeaways clave}

\begin{enumerate}
\item Los datos sintéticos no sustituyen a los datos reales: parametrizan y llevan a gran escala los problemas de escasez del mundo real, y se calibran mediante retroalimentación real.
\item La inteligencia corporizada depende más que la conducción autónoma de la interacción física, por lo que la unidad de datos pasa de frame a interaction.
\item Una buena simulación no es la que tiene una imagen realista, sino la que sirve para entrenar modelos, ejecutar RL en paralelo y llevar los resultados a robots reales.
\item Las empresas de datos se parecerán cada vez más a empresas de infraestructura central de IA, porque los buenos datos requieren talento, procesos, control de calidad, escenarios y conocimiento del cliente.
\item La inteligencia corporizada sigue en la etapa de GPT-1; el verdadero scaling law moment depende de si la capacidad en tareas reales mejora de forma estable con más datos, cómputo y entornos.
\end{enumerate}

\subsection{Lecciones para el flujo de trabajo}

Si se convierte EP109 en una checklist interna para un equipo, puede hacerse una autoevaluación con cuatro preguntas. Primera: ¿tenemos métricas claras del modelo objetivo, o solo estamos generando datos? Segunda: ¿nuestros activos de simulación incluyen parámetros físicos con los que se puede interactuar, y no solo la apariencia visual? Tercera: ¿tenemos despliegues reales o retroalimentación de clientes que le indiquen al equipo de datos dónde está el gap? Cuarta: ¿nuestro reward model se basa en el progress, la recompra y la capacidad en tareas reales, o en la demo, el financiamiento y la narrativa?

Esta es también la lección práctica que más vale conservar de este episodio: quien trabaja en datos sintéticos y simulación debe partir del usuario final y razonar hacia atrás. El sistema de producción de datos debe evolucionar en torno a estas preguntas: qué interfaces necesita el equipo de algoritmos de robótica, qué environment necesita el RL, qué escenarios necesita el cliente y qué fallas revela el despliegue real.

\subsection{Lecturas adicionales}

\begin{itemize}
\item Si le interesan los modelos del mundo para robots, la transferencia entre distintos cuerpos robóticos y Gemini Robotics, puede compararse con la entrevista a Tan Jie (谭捷), de DeepMind, en EP121.
\item Si le interesan el robot completo, la cadena de suministro y la Data Recipe, puede compararse con la entrevista a Gao Jiyang (高继扬), de Xinghaitu, en EP132.
\item Si le interesan el valor de los datos para la IA, Meta/Scale y la estrategia de las empresas de datos, puede compararse con las entrevistas de Zhang Xiaojun sobre datos para la IA.
\end{itemize}

\end{document}
